\documentclass[12pt, a4paper]{ctexart}
\usepackage{geometry}
\geometry{left=2.5cm, right=2.5cm, top=3cm, bottom=3cm}
\usepackage{amsmath}
\usepackage{circuitikz}
\usepackage{float}
\usepackage{listings}
\usetikzlibrary{arrows, shapes.gates.logic.US}
\ctikzset{
    logic ports=ieee,
    logic ports/scale=0.8,
    multipoles/flipflop/pin spacing=0.4
}
\input{risc-v.tex}
\lstset{
    basicstyle=\ttfamily\small,
    frame=single,
    framesep=5pt,
    xleftmargin=10pt,
    xrightmargin=10pt,
    breaklines=true
}

\begin{document}

\section*{第7章习题}

\subsection*{2.}

\subsubsection*{(4)}

\begin{itemize}
    \item 在寄存器中：寄存器直接寻址
    \item 在存储器中：寄存器间接寻址，直接寻址，间接寻址，基址寻址，变址寻址，相对寻址
\end{itemize}

\subsubsection*{(9)}

\begin{itemize}
    \item 转移跳转指令不会返回执行，调用指令会返回继续执行。
    \item 一般不需要，因为 CPU 在跳转之前会自动将下一条指令的地址压入栈中，
    或者放入特定的寄存器 （LR）中，返回时 CPU 会直接读取存在特定位置的地址，
    不需要提供地址码字段。不过如果是较早期的 CPU，可能需要。
\end{itemize}

\subsection*{3.}

假定题目中的 $PC=258$ 为取指前的 $PC$，则取指后 $PC=260$，位移量为 $-40$，
补码为 \texttt{D8H}，所以该转移指令第二字节的内容应为 \texttt{D8H}。

\subsection*{4.}

\subsubsection*{(1)}

$EA=A+(I)=\mathtt{1200H}+\mathtt{FCH}=\mathtt{12FCH}$

\subsubsection*{(2)}

$EA=(A)=\mathtt{12FCH}$

\subsubsection*{(3)}

$EA=(R_8)=\mathtt{1200H}$

\subsection*{6.}

$k_1=1024-\dfrac{k_0}{64}-64k_2$

\subsection*{7.}

\begin{table}[H]
    \centering
    \begin{tabular}{c|c|c|c|c|c|c|c|}
        \multicolumn{2}{c}{} & \multicolumn{1}{c}{Ms[2]} &
        \multicolumn{1}{c}{} & \multicolumn{1}{c}{Md[2]} \\
        \cline{2-6}
        RR & OP[6] & 01 & Rs[3] & 01 & Rd[3] \\ \cline{2-7}
        RI & OP[6] & 00 & 000   & 01 & Rd[3] & Imm[16] \\ \cline{2-7}
        RS & OP[6] & 01 & Rs[3] & 10 & Rd[3] \\ \cline{2-7}
        RX & OP[6] & 01 & Rs[3] & 11 & Rx[3] & Ad[16]  \\ \cline{2-8}
        XI & OP[6] & 00 & 000   & 11 & Rx[3] & Imm[16] & Ad[16] \\ \cline{2-8}
        SI & OP[6] & 00 & 000   & 10 & Rd[3] & Imm[16] \\ \cline{2-7}
        SS & OP[6] & 10 & Rs[3] & 10 & Rd[3] \\ \cline{2-6}
    \end{tabular}
\end{table}

其中 $OP$ 表示 6 位操作码；$M_S$ 和 $M_D$ 分别表示源操作数和目的操作数的 2 位寻址方式特征码；
$R_S$ 和 $R_D$ 分别表示源操作数和目的操作数所使用的 3 位通用寄存器编号
（当寻址方式为变址时，$R_X$ 特指用作变址寄存器的 3 位通用寄存器编号）；
$IMM$ 表示 16 位立即数，$A_D$ 表示 16 位目的操作数形式地址。

\subsection*{8.}

\subsubsection*{(1)}

16 条指令，8 个通用寄存器，16 位 MAR 和 16 位 MDR。

\subsubsection*{(2)}

\texttt{0000H} 至 \texttt{FFFFH}。

\subsubsection*{(3)}

\texttt{add (R4), (R5)+} 对应 \texttt{0010 001 100 010 101B}。

执行后 \texttt{R5} 的内容变为 \texttt{5679H}，\texttt{5678H} 的内容变为 \texttt{68ACH}。

\subsection*{9.}

假设 \texttt{s0 = A\_upper20, s1 = A\_lower12}，
输出 \texttt{A\_upper20\_adjusted} 在 \texttt{a0}，
则用以下方式可以得到

\subsubsection*{(1)}

\begin{lstlisting}[language={[RISC-V]Assembler}]
srai    a0, s1, 11
xor     a0, s0, a0
\end{lstlisting}

\subsubsection*{(2)}

\begin{lstlisting}[language={[RISC-V]Assembler}]
srli    a0, s1, 11
add     a0, s0, a0
\end{lstlisting}

\subsection*{10.}

当 \texttt{offset[11]} 为 0 时应为 \texttt{offset[31:12] + 0}，
当 \texttt{offset[11]} 为 1 时应为 \texttt{offset[31:12] + 1}，
所以高 20 位可以表示为 \texttt{offset[31:12] + offset[11]}。

\subsection*{11.}

\begin{lstlisting}[language={[RISC-V]Assembler}]
slli    t1, t0, 3
sub     t1, t1, t0
\end{lstlisting}

\subsection*{13.}

\begin{lstlisting}[language={[RISC-V]Assembler}]
loop:
  add   t0, zero, zero    # 初始化寄存器：t0 = 0 + 0
  beq   a1, zero, finish  # 条件跳转：如果 a1 == 0，跳转到 finish
  add   t0, t0, a0        # 累加：t0 = t0 + a0
  addi  a1, a1, -1        # 计数器递减：a1 = a1 - 1
  j     loop              # 无条件跳转：回到 loop 标签，继续循环

finish:
  addi  t0, t0, 100       # 加偏置：t0 = t0 + 100
  add   a0, t0, zero      # 返回值赋值：a0 = t0 + 0
\end{lstlisting}

\subsection*{15.}

\subsubsection*{(1)}

\begin{lstlisting}[language={[RISC-V]Assembler}]
addi    t1, t0, 31
\end{lstlisting}

\subsubsection*{(2)}

\begin{lstlisting}[language={[RISC-V]Assembler}]
slli    t1, t0, 16
srli    t1, t1, 16
\end{lstlisting}

\subsection*{17.}

\texttt{beq} 含义：Branch if EQual 相等则跳转

问题：offset 只有 12 位，\texttt{here} 和 \texttt{there} 之间的指令数量可能超过 offset

修改：

\begin{lstlisting}
here:
  bne   t0, t2, skip
  j     there
skip:
  ......
there:
  addi  t1, a0, 4
\end{lstlisting}

\end{document}